% docs/documento_oficial/abstract.tex
\section*{Abstract}
\addcontentsline{toc}{section}{Abstract}

In the Meta department (Colombia), environmental conditions accelerate concrete degradation and increase maintenance costs. Given the lack of local tools to support mix design and admixture selection, this project develops DIAPCE: a cross-platform Flutter application aimed exclusively at students in the research seedbed and participants in materials laboratory practices at the Universidad Cooperativa de Colombia (UCC), Meta campus. DIAPCE is not an industrial or general-purpose tool; its purpose is purely educational and institutional. Based on a regional dataset and inputs such as target strength, temperature, humidity, and exposure conditions, the system suggests a base mix design and classifies the use of admixtures (P0: no admixture; P1-P3: waterproofing at 2, 3, and 4\,\%; PP1-PP3: plasticizer at 0.2, 0.4, and 0.6\,\%), with projections at 7, 14, and 28 days, in accordance with ACI\,211.1 and related standards. The solution works offline using SQLite and prioritizes the traceability of academic tests and exercises (recording of inputs, results, and model version), facilitating reproducible and comparable practices in the classroom and laboratory. DIAPCE seeks to strengthen skills in mix design, results analysis, and data-recording culture among UCC students, potentially contributing to a more efficient use of materials within the context of their educational practices.

\vspace{0.5em}
\noindent\textbf{Keywords:} structural concrete, mix design, admixtures (P0-PP3: waterproofing 0-4\,\% and plasticizer 0.2-0.6\,\%), ACI\,211.1, Flutter, SQLite, Meta (Colombia), 7/14/28 days, laboratories, research seedbeds, teaching.
